\chapter{What drives the performance}
\label{ch:drivers}
% This chapter is what makes the report more than "my model got a higher
% number". You tested your own design and reported what you found.

\section{Ablation on LOFAR}
% Method [PART 17, 18]: experiments/models_ablation.py makes each component
% a switch; experiments/lofar_hybrid.py --variant trains every variant
% through the identical pipeline. Control: the switchable model with
% everything on scored 0.6624 vs the real model's 0.6603 +/- 0.0040.
% Seed 0 table [PART 18.2]:
%   hybrid_full 0.6624 | no_strip 0.6637 | no_eca 0.6557 | no_res 0.6569 |
%   plain_unet 0.6605 | no_groupnorm 0.6289
% 3-seed table [PART 20]:
%   full hybrid    0.6603 +/- 0.0040
%   plain_unet     0.6585 +/- 0.0047   (strip, ECA, residual all removed)
%   no_groupnorm   0.6495 +/- 0.0178   (also no normalisation)
%   Welch t: full vs plain p = 0.64; plain vs no_groupnorm p = 0.48.
% Conclusion: the three added components contribute nothing measurable.
% GroupNorm's effect is not established either -- seed 0 suggested a large
% effect that vanished at 3 seeds. Without GroupNorm, training is less
% stable (sd 0.0178 vs ~0.004).

\section{Do strip convolutions help faint RFI?}
% - The design claim: strips integrate along lines so faint RFI becomes
%   detectable. Overall F1 can't test this (faint RFI is a small share),
%   so you measured recall per brightness tier.
% - First attempt was misleading: each model at its own threshold showed
%   ~+0.02 for strips in the dim tiers -- but hybrid_full's threshold was
%   0.296 vs no_strip's 0.434, and a lower threshold raises recall
%   everywhere. [PART 18.5]
% - Re-done at matched budget (every model flags exactly 199,951 pixels):
%   bright -0.0053 | moderate +0.0029 | faint -0.0008 | invisible +0.0008.
%   No effect anywhere. [analysis/ablation_faint_recall.py]
% - Worth a sentence on the lesson: compare models at the same operating
%   point.

\section{Testing a CRF refinement}
% [PART 16] Chen et al.'s CRF idea, built as a differentiable layer on top
% of the trained network (11 extra parameters).
%  - Chen et al.'s own parameters: 0.6592 -> 0.4884 (destroys thin RFI).
%  - Best hand-tuned: +0.0012. Trained end-to-end: +0.0012 (0.29 seed sd).
%  - Given freedom, the CRF shrank its own weights toward zero.
%  - Why: the network already uses neighbourhood context, and a CRF can only
%    redistribute evidence that is already in the prediction.

\section{Where the remaining errors are}
% [PART 15] Recall by tier for the best model:
%   bright 0.937 | moderate 0.391 | faint 0.268 | invisible 0.254
% - Bright RFI is solved; the dim half is caught only a quarter to a third
%   of the time.
% - Largest possible gain: removing false positives (+0.114 upper bound);
%   perfectly catching the invisible tier (+0.098).
% - And recall the ceiling: training labels are ~44% wrong.

\section{So what does explain the gain over \tfunet?}
% [PART 20.2] Not the three components and not clearly GroupNorm. Still
% untested and confounded: Dice loss, raw-logit output (no ReLU), same
% padding, depth/width, batch size, learning rate, framework.
% Be direct: this is not yet known, and say which experiment comes next
% (start with the Dice loss).
